% Droite d'étalonnage des ions cuivre(II) à 800 nm (exercice 4, chapitre 2)
\begin{tikzpicture}
\begin{axis}[otMesures,otAxesBout,otLecture,width=110mm,height=55mm,xmin=0,xmax=90,ymin=0,ymax=1.1,
  xtick={0,10,...,80},ytick={0,0.2,...,1.0},
  yticklabel style={/pgf/number format/fixed,/pgf/number format/fixed zerofill,/pgf/number format/precision=1},
  xlabel={$c$ (\unit{mmol.L^{-1}})},ylabel={$A$ à \qty{800}{nm}}]
\addplot[only marks,mark=*,mark size=1.7pt,otInk] coordinates {(10,0.12) (20,0.24) (40,0.48) (60,0.72) (80,0.96)};
\solution{\addplot[otGreenDark,line width=0.8pt,domain=0:88,samples=2] {0.012*x};
  \draw[otOrange,line width=0.6pt,dash pattern=on 2pt off 1.5pt] (axis cs:0,0.54) -- (axis cs:45,0.54) -- (axis cs:45,0);
  \fill[otOrange] (axis cs:45,0.54) circle (1.6pt);
  \node[pcLegende,anchor=north west,text=otOrange!75!black,fill=white,inner sep=1pt] at (axis cs:46,0.3) {$c=\qty{45}{mmol.L^{-1}}$};}
\end{axis}
\end{tikzpicture}
\par\smallskip\legende{Absorbance de solutions de sulfate de cuivre de concentrations connues (cuve de \qty{1,00}{cm}).}
